
%\begin{Verbatim}[frame=single,framesep=6.5mm,label=\fbox{Example file 1 },numbers=left,numbersep=-10pt,xrightmargin=5cm]
%\end{Verbatim}

\chapter{Examples \& Tutorials}
\label{sec:Examples Tutorials}
% Andreas / Simon

% Input und Output kurz erläutern.

This part of the manual deals with example applications, where the use of the provided software will be demonstrated. The examples presented in this section were chosen such that a wide class of applications is covered. Each problem is briefly described and the appropriate input files are given. also includes some output files. 


\section{Lennard-Jones fluid}
\label{sec:Lennard-Jones Fluid}

An example for a complete *.par file is given below for a MD simulation of the Lennard-Jones fluid in the \textit{NVT} ensemble, where the chemical potential is calculated by the Widoms's test molecule method. In addition, the calculation of the transport properties is enabled.



Parameters and options specified in the *.pm file for the Lennard-Jones fluid:\\\\

\begin{Verbatim}[frame=single,framesep=6.5mm,label=\fbox{Lennard-Jones fluid},numbers=left,numbersep=-10pt,xrightmargin=5cm]
# MS2 Parameterfile created with ms2par
# General simulation parameters
Units         =       Reduced
LengthUnit    =       1.0
EnergyUnit    =       1.0
MassUnit      =       1.0 
Simulation    =       MD
Integrator    =       Gear
TimeStep      =       0.0010
Ensemble      =       NVT
MCORSteps     =       0
NVTSteps      =  300000
RunSteps      = 1000000
ResultFreq    =    1000
ErrorsFreq    =    5000
VisualFreq    =       0
CutoffMode    =       COM
NEnsembles    =       1
CorrfunMode   =       yes
#
# Ensemble 1
#
Temperature   =       1.10
Density       =       0.66
OptPressure   =       Yes
NParticles    =    1372
NComponents   =       1
Corrlength    =    6000
SpanCorrfun   =     300
ViewCorrfun   =     100
ResultFreqCF  =       1

PotModel      =       lj.pm
MoleFract     =       1.0
ChemPotMethod =       Widom
NTest         =    3000

Cutoff        =       6.3
Epsilon       =       1.0E10
\end{Verbatim}

\vspace{1cm}






\section{Vapour-Liquid Equilibrium of a \ce{O2} -- Aceton mixture}
\label{sec:Vapour-Liquid Equilibrium of a O2 - Aceton mixture}

\begin{Verbatim}[frame=single,framesep=6.5mm,label=\fbox{VLE \ce{O2} -- Aceton mixture},numbers=left,numbersep=-10pt,xrightmargin=5cm]
#
# MS2 Parameterfile created with ms2par
#

#
# General simulation parameters
#
Units       =       SI
LengthUnit  =       3.5
EnergyUnit  =       100.0
MassUnit    =       40.0
Simulation    =      MC
Acceptance    =      0.5
Ensemble      =      GE
NVTSteps      =   10000
mueVTSteps    =   25000
RunSteps      =  200000
ResultFreq    =     100
ErrorsFreq    =    1000
VisualFreq    =       0
CutoffMode  =       COM
NEnsembles    =       1

#results from liquid run
Temperature = 400.00   
Pressure0     =       9.000
LiqDensity    =      11.5790
VarDensity    =       0.0051
LiqEnthalpy   =  -23452.3
VarEnthalpy   =       8.4
LiqBetaT      =       0.004760
VarBetaT      =       0.000252
LiqdHdP       =      -0.001457
VardHdP       =       0.004658

#parameters for vapor run
Density       =       1.000
OptPressure   =     Yes
NParticles    =     500
NComponents   =       2

PotModel      =      O2.pm
MolarFract    =       0.5000
LiqMolarFract =       0.0799
ChemPot       =      -2.849
VarChemPot    =       0.003
PartMolVol    =       3.4185
VarPartMolVol =       0.1899


PotModel      =  aceton.pm
MolarFract    =       0.5000
LiqMolarFract =       0.9201
ChemPot       =      -5.246
VarChemPot    =       0.011
PartMolVol    =       3.5289
VarPartMolVol =       0.3598

eta         = 1.0
xi          = 0.905
Cutoff      = 4.0
Epsilon     = 1.0E10
\end{Verbatim}



\section{calculating configurational properties}
\label{sec:calculating configurational properties}





\section{calculating transport properties}
\label{sec:calculating transport properties}






\section{chemical potential by ThermoInt}
\label{sec:chemical potential by ThermoInt}




\section{\ce{NaCl}-Water Solution \& osmotic pressure}
\label{sec:NaCl-Water Solution osmotic pressure}





\section{multi-Ensemble method}
\label{sec:multi-Ensemble method}




\section{humid air ensemble}
\label{sec:humid air ensemble}















